\chapter{Fundamentação teórica}\label{cap:fundamentacao}

Este capítulo apresenta os conceitos necessários à compreensão do trabalho:
o fluxo de projeto de uma PCB, os mecanismos de automação disponíveis em
ferramentas de EDA, a estrutura do KiCad, os fundamentos de agentes baseados
em LLMs e o \emph{Model Context Protocol}.

\section{O fluxo de projeto de uma PCB}

O desenvolvimento de uma placa de circuito impresso pode ser descrito como uma
sequência de etapas com critérios de verificação próprios
\citep{kicaddocs,horowitz2015}:

\begin{enumerate}[leftmargin=1.4em,itemsep=0.2em]
  \item \textbf{Especificação funcional}: definição dos blocos elétricos,
        tensões, correntes e interfaces do sistema;
  \item \textbf{Captura esquemática}: representação da conectividade por
        símbolos e fios, com anotações de valores e encapsulamentos;
  \item \textbf{Atribuição de \emph{footprints}}: associação de cada símbolo
        ao desenho físico correspondente (dimensões, pads, \emph{courtyard});
  \item \textbf{Layout}: posicionamento dos componentes e definição da pilha
        de camadas, do contorno da placa e das zonas de cobre;
  \item \textbf{Roteamento}: traçado das trilhas e posicionamento de vias,
        respeitando regras elétricas e de fabricação;
  \item \textbf{Verificação}: checagem de regras de projeto (DRC) para
        geometria e de regras elétricas (ERC) para o esquemático, incluindo a
        paridade entre placa e esquemático;
  \item \textbf{Preparação para fabricação}: geração de Gerber, furação,
        máscara, lista de materiais e arquivos de posição de montagem.
\end{enumerate}

As verificações das etapas 6 e 7 são especialmente importantes para a
automação, pois produzem resultados estruturados (listas de violações por
tipo e severidade) que podem ser processados por ferramentas externas. Em
contrapartida, operações das etapas 2 a 5 são destrutivas do ponto de vista do
arquivo de projeto: alteram geometria e conectividade, exigindo que qualquer
automação controle salvamento, recarga e estado de sessão.

\section{Automação em ferramentas de EDA}

Ferramentas de EDA oferecem, tipicamente, três níveis de automação:

\begin{itemize}[leftmargin=1.4em,itemsep=0.2em]
  \item \textbf{Interfaces de programação em processo} (\emph{in-process
        scripting}): o usuário executa código dentro do processo da ferramenta,
        com acesso direto ao modelo de dados. No KiCad, essa via é a API
        Python \texttt{pcbnew}, gerada por \emph{bindings} SWIG
        \citep{swig}.
  \item \textbf{Interfaces de comunicação entre processos}: a ferramenta expõe
        uma API de serviço com a qual aplicações externas conversam em tempo
        real. No KiCad, essa via é a API IPC \citep{kicaddev}.
  \item \textbf{Utilitários de linha de comando}: operações de verificação e
        exportação executadas sem interface gráfica. No KiCad, essa via é o
        \texttt{kicad-cli}, com subcomandos para DRC, ERC, exportações Gerber,
        furação, PDF, SVG, STEP, entre outros \citep{kicadcli}.
\end{itemize}

Cada nível tem vantagens e limites: a programação em processo tem acesso
completo, mas exige um ambiente Python compatível com a instalação da
ferramenta; a API IPC é adequada para sincronização em tempo real, porém
depende de um servidor habilitado pelo usuário; a linha de comando é robusta e
isolada, mas não mantém estado entre execuções. Uma integração madura
beneficia-se dos três, escolhendo o mecanismo conforme a operação.

\section{O KiCad}

O KiCad é uma suíte de EDA de código aberto que integra captura esquemática,
editor de \emph{footprints}, editor de símbolos, editor de placas, visualizador
3D e utilitários de fabricação \citep{kicaddocs}. Os arquivos de projeto são
armazenados em formato textual de expressões-S (arquivos
\texttt{.kicad\_pro}, \texttt{.kicad\_sch}, \texttt{.kicad\_pcb} e bibliotecas
\texttt{.kicad\_sym}/\texttt{.kicad\_mod}), o que favorece controle de versão
e inspeção por ferramentas externas.

Duas características são relevantes para este trabalho. A primeira é a
\textbf{estabilidade da API de automação}: a API \texttt{pcbnew} cobre placas,
enquanto a manipulação de esquemáticos é feita por ferramentas próprias da
suíte; por isso, integrações de esquemático frequentemente combinam geração e
leitura de expressões-S com validação pelo \texttt{kicad-cli}. A segunda é a
\textbf{versatilidade de verificação}: os relatórios de DRC e ERC do
\texttt{kicad-cli}, emitidos em formato JSON, permitem automação de
qualidade --- exatamente os artefatos utilizados na avaliação deste trabalho
\citep{kicadcli}.

\section{Modelos de linguagem e agentes}

Modelos de linguagem modernos são capazes de manter contexto sobre
documentação e código, planejar sequências de ações e produzir chamadas
estruturadas a ferramentas. A literatura documenta três marcos relevantes para
esta dissertação:

\begin{itemize}[leftmargin=1.4em,itemsep=0.2em]
  \item Yao et al.~\cite{yao2023react} demonstram que intercalar raciocínio e ação, com
        observações do ambiente realimentando o ciclo, melhora o desempenho em
        tarefas que exigem decisões encadeadas;
  \item Schick et al.~\cite{schick2023toolformer} e Patil
        et al.~\cite{patil2023gorilla} evidenciam que
        modelos podem aprender a usar APIs externas de forma autônoma quando
        as chamadas lhes são apresentadas de forma estruturada;
  \item Yang et al.~\cite{yang2024sweagent} mostram que a interface entre agente e
        computador (ACI) é determinante: comandos claros, saídas concisas e
        ferramentas descobríveis elevam o desempenho do agente.
\end{itemize}

Esses resultados orientam decisões concretas adotadas no servidor estudado:
exposição individual de cada ferramenta, com descrição e esquema próprios;
lista de ferramentas essenciais para ordenar a descoberta; respostas
estruturadas e concisas; e separação clara entre operações de leitura e de
escrita.

\section{O Model Context Protocol}

\subsection{Origem e motivação}

O MCP foi anunciado pela Anthropic em novembro de 2024 como um padrão aberto
para conectar assistentes de IA a sistemas externos \citep{anthropic2024mcp}.
A motivação central é reduzir o problema $N \times M$: sem um protocolo comum,
cada combinação de assistente e ferramenta requer uma integração dedicada;
com o protocolo, ferramentas publicam um servidor e assistentes implementam um
cliente.

\subsection{Arquitetura}

A arquitetura do MCP define três papéis \citep{mcp2025spec}:

\begin{itemize}[leftmargin=1.4em,itemsep=0.2em]
  \item \textbf{Anfitrião} (\emph{host}): aplicação que hospeda o assistente e
        coordena um ou mais clientes;
  \item \textbf{Cliente}: componente que mantém a conexão com um servidor e
        encaminha requisições;
  \item \textbf{Servidor}: programa que expõe capacidades (ferramentas,
        recursos, prompts) e responde às requisições.
\end{itemize}

A comunicação usa JSON-RPC 2.0, com transporte por entrada/saída padrão
(\emph{stdio}) ou por HTTP com \emph{streaming}. O ciclo de vida inclui
inicialização com negociação de capacidades e notificações de mudança. A
Figura~\ref{fig:mcp-geral} ilustra a organização geral aplicada a este
trabalho.

\begin{figure}[htbp]
\centering
\begin{tikzpicture}[
  font=\footnotesize,
  box/.style={draw, rounded corners=2pt, align=center, inner sep=5pt},
  arr/.style={-{Stealth[length=2mm]}, thick},
  node distance=8mm and 10mm
]
\node[box, fill=blue!6] (host) {Aplicação anfitriã\\(editor de código, agente)};
\node[box, fill=blue!12, right=of host] (srv) {Servidor MCP\\\texttt{mcp-kicad-vscode}\\(233 ferramentas, 23 recursos)};
\node[box, fill=green!10, right=of srv] (kicad) {KiCad 9\\\texttt{pcbnew} (SWIG), API IPC,\\\texttt{kicad-cli}};
\draw[arr] (host) -- node[above, font=\scriptsize]{MCP} (srv);
\draw[arr] (srv) -- node[above, font=\scriptsize]{comandos} (kicad);
\draw[arr] (kicad.west) -- ++(0,-1.6) -| node[below, font=\scriptsize, pos=0.25]{resultados} (srv.south);
\end{tikzpicture}
\caption{Organização geral: o assistente conecta-se ao servidor MCP, que opera
o KiCad e devolve resultados estruturados.}
\label{fig:mcp-geral}
\end{figure}

\subsection{Primitivas}

O protocolo define três primitivas de servidor, resumidas no
Quadro~\ref{qua:primitivas} \citep{mcp2025docs}.

\begin{quadro}[htbp]
\centering
\small
\caption{Primitivas do MCP.}
\label{qua:primitivas}
\begin{tabular}{@{}p{2.4cm}p{3.4cm}p{8.4cm}@{}}
\toprule
\textbf{Primitiva} & \textbf{Direção} & \textbf{Função} \\
\midrule
\emph{Tools} & modelo $\rightarrow$ servidor & Operações executáveis com
efeitos, descritas por esquema de parâmetros; exigem confirmação do usuário
quando destrutivas. \\
\emph{Resources} & servidor $\rightarrow$ modelo & Dados de contexto somente
leitura, identificados por URI (estado do projeto, listas, metadados). \\
\emph{Prompts} & servidor $\rightarrow$ modelo & Modelos de interação
reutilizáveis para tarefas recorrentes. \\
\bottomrule
\end{tabular}
\end{quadro}

Além das primitivas, a especificação define elementos transversais relevantes
para integrações robustas: \emph{timeouts}, cancelamento de requisições,
notificações de progresso e limites de tamanho de mensagem
\citep{mcp2025spec}.

\subsection{Adoção e relevância}

Após o anúncio, o protocolo passou a ser implementado por diferentes
fornecedores e comunidades, e servidores MCP foram publicados para domínios
distintos --- de edição de código a sistemas de arquivos e bancos de dados.
Para o presente trabalho, a adoção importa por dois motivos: valida a aposta
de padronização e amplia o número de assistentes capazes de consumir o
servidor sem código adicional.

\section{Síntese do capítulo}

A fundamentação apresentada permite situar o problema: o fluxo de PCB combina
operações destrutivas e verificações estruturadas; o KiCad oferece três
mecanismos de automação complementares; os agentes baseados em LLMs exigem
interfaces desenhadas para sua forma de operação; e o MCP fornece o contrato
padronizado que faltava. O próximo capítulo examina como trabalhos
relacionados tratam essa combinação.
